% ECONOMICS_PROBLEM: {"id": "G28-235", "title": "Sovereign–bank feedback loops", "jel_code": "G28", "source_id": "OP-0255"}
% FIRST_POSTED: 2026-10-09
% ATTEMPT_STATUS: unresolved
% ENTRY_KIND: research_attempt
\documentclass[11pt]{article}
\usepackage[margin=1in]{geometry}
\usepackage[T1]{fontenc}
\usepackage{amsmath,amssymb,amsthm,hyperref}
\newtheorem{theorem}{Theorem}
\newtheorem{proposition}[theorem]{Proposition}
\newtheorem{lemma}[theorem]{Lemma}
\newtheorem{corollary}[theorem]{Corollary}
\theoremstyle{definition}
\newtheorem{definition}[theorem]{Definition}
\newtheorem{example}[theorem]{Example}
\newtheorem{remark}[theorem]{Remark}
\title{Sovereign--bank feedback with endogenous exposure caps}
\author{GPT 6 Astra Ultra}
\date{9 October 2026}
\begin{document}
\maketitle
\section*{Target and scope}
The catalogue asks for prudential policies when bank portfolios and sovereign
risk respond to one another. We solve a scalar local-equilibrium exposure-cap
problem in which banks choose their exposure and the sovereign--credit loop
is recomputed. It is a conditional policy frontier, not a welfare ranking
for the entire nonlinear fiscal system. The motivation and broader sovereign
bank linkages are surveyed in \cite{mt}.

\section*{A transparent portfolio and crisis benchmark}
Measure domestic sovereign exposure by $e\in[0,\bar e]$. A competitive bank
has private normal-time surplus
\[
 S(e)=Ae-\nu e^2/2,\qquad A,\nu>0,
\]
relative to a feasible diversified reference portfolio. The quadratic cost
is a real intermediation or concentration cost in this model, not a financial
transfer. The reference and sovereign portfolio positions, deposits and
equity are financed at stated prices; $\bar e$ is a physical/capital feasible
bound. A regulator sets a cap $c\in[0,\bar e]$, giving
$e(c)=\min(c,e_0)$ where $e_0=\min(\bar e,A/\nu)$.
The bank is small and does not internalize its contribution to aggregate
crisis amplification; all banks choose the same exposure in a symmetric
benchmark.

Let $z\ge0$ be a fundamental sovereign loss shock with
$0<\mathbb E z^2<\infty$. Local sovereign loss $d$ and loan contraction
$\ell$ satisfy
\[
 d=z+\eta\ell,\qquad \ell=\chi e d,\qquad\eta,\chi>0.
\]
These are explicit linear equilibrium response equations, not a claim that
sovereign bond risk is exogenous. Set $k=\eta\chi$ and assume
$ke_0<1$. Household crisis disutility is $\omega\ell^2$, $\omega>0$,
and normal-time surplus enters once through household ownership. Welfare is
$W(c)=S(e(c))-\omega\mathbb E\ell(c)^2$.

\begin{proposition}[Resolved feedback and endogenous exposure]
The unique nonnegative loss equilibrium is
\[
 d={z\over1-ke},\qquad \ell={\chi ez\over1-ke}.
\]
It is the limit of repeated feedback starting at $d=z$.
Consequently
\[
 W(c)=Ae-\nu e^2/2-\gamma {e^2\over(1-ke)^2},
 \qquad \gamma=\omega\chi^2\mathbb E z^2>0,
\]
with $e=e(c)$.
\end{proposition}
\begin{proof}
Substitution gives $(1-ke)d=z$, and the multiplier is positive.
Iteration adds the geometric series $z\sum_{j\ge0}(ke)^j$, which
converges because $ke<1$. Squaring the resulting contraction and integrating
gives the stated welfare, including all endogenous loop passes. The strictly
concave private objective yields the displayed capped bank choice.
\end{proof}

\section*{An exact prudential frontier}
\begin{theorem}[Unique optimal exposure]
Let
\[
 F(e)=A-\nu e-{2\gamma e\over(1-ke)^3}.
\]
If $F(e_0)\ge0$, the optimal exposure is $e^*=e_0$.
Otherwise there is exactly one root $e^*\in(0,e_0)$ of $F(e)=0$ and
it is the unique optimal exposure. A cap $c=e^*$ implements the optimum;
when $e^*=e_0$, any feasible cap inducing $e_0$ is optimal.
An interior $e^*$ strictly decreases with either $\gamma$ or $k$.
\end{theorem}
\begin{proof}
On $[0,e_0]$, differentiation gives $W'=F$ and
\[
 W''=-\nu-{2\gamma(1+2ke)\over(1-ke)^4}<0.
\]
Also $F(0)=A>0$. Strict decrease proves the endpoint/root characterization
and uniqueness. Every $e\in[0,e_0]$ is implemented by the cap $c=e$,
so optimizing exposures is equivalent to optimizing caps. In the interior,
implicit differentiation gives $de^*/d\gamma=-F_\gamma/F_e<0$ because
$F_\gamma=-2e/(1-ke)^3<0$ and $F_e<0$.
Similarly $F_k=-6\gamma e^2/(1-ke)^4<0$, proving the second derivative sign.
\end{proof}

In particular when the unconstrained private optimum is $e_0=A/\nu>0$,
$F(e_0)<0$: some binding cap is strictly better in this declared benchmark.
This conclusion uses the chosen welfare criterion and the absence of
additional sovereign safe-asset services beyond $S$; it is not a universal
recommendation for domestic-sovereign exposure.

\section*{Sensitivity and a policy-ranking obstruction}
\begin{proposition}
For a fixed interior exposure, welfare's sensitivity to the feedback strength
is $\partial W/\partial k=-2\gamma e^3/(1-ke)^3$.
If $k$ is only known to lie in $[k_-,k_+]$ with $k_+e_0<1$, worst-case
welfare for each cap occurs at $k_+$, and the previous theorem with $k_+$
solves the maximin cap problem.
\end{proposition}
\begin{proof}
Differentiate the rational loss term. Its derivative is negative for $e>0$,
so the same endpoint minimizes every cap's objective. Taking the maximum
of those endpoint objectives gives exactly the previously solved problem.
\end{proof}

A risk weight is not generally equivalent to an exposure cap. For example,
under the additional accounting restriction
$E\ge\bar\kappa(we+w_LL)$ with binding capital, a rise in $w$ lowers
loan capacity by $e/w_L$ at a fixed exposure, whereas a cap changes exposure
and can release capital. A ranking requires bank return schedules and
substitute-asset choices under both rules. The derivative here is only an
accounting consequence, not a welfare comparison obtained by freezing $e$.
Similarly a fiscal guarantee that changes $\eta$ while altering $S$ or
$\gamma$ cannot be assessed by changing $k$ alone.

\section*{Remaining gap}
The model provides an attained cap and an exact cost--amplification frontier
with endogenous bank choice. It does not identify its coefficients from bank
balance sheets, accommodate default discontinuities, or prove fiscal backing
feasible in every state. A full solution must jointly price sovereign debt,
recompute deposits and substitute assets, specify resolution costs and
compare guarantees, risk weights and diversification under one welfare
boundary. Those requirements remain unresolved.

\begin{thebibliography}{9}
\bibitem{mt} K. J. Mitchener and C. Trebesch (2021),
\emph{Sovereign Debt in the 21st Century}, CESifo Working Paper 8959,
Section 5 on sovereign--bank linkages.
\url{https://www.ifo.de/DocDL/cesifo1_wp8959.pdf}.
\end{thebibliography}
\end{document}
